% =====================================================================
%  c01.tex —— §2 A∞ 范畴（原书 p.44–46；PDF p.3（§2 起）–p.5（§3 标题前））
%  编号公式：式 (2)（A∞ 关系式，原书 p.44）、式 (3)（A∞ 函子的多重线性映射族，原书 p.45）
%  定理类环境：定义 2.1、定义 2.2、定理 2.3
%  说明：§3 标题及其后的式 (4) 属 c02，本块不含。
% =====================================================================

% ===== 原书 p.44 =====
\section{\texorpdfstring{$\ainf$范畴}{A∞ 范畴}}
\label{sec:c01}

$\ainf$范畴是把$\ainf$代数[21, 22]的元素提升为范畴的态射而得到的东西[3]。
也就是说，$\ainf$范畴$\mathcal{C}$是指下述数据。对象集合
$\Ob(\mathcal{C})=\{X,Y,\dots\}$；对任意两个对象$X,Y\in\Ob(\mathcal{C})$，有
$\mathbb{Z}$分次向量空间
$\mathcal{C}(X,Y)=\bigoplus_{r\in\mathbb{Z}}\mathcal{C}^{r}(X,Y)$；
此外还有一族次数为$(2-n)$的多重线性映射（multilinear map）
\[
  m_{n}:\mathcal{C}(X_{1},X_{2})\otimes\cdots\otimes\mathcal{C}(X_{n},X_{n+1})
  \longrightarrow\mathcal{C}(X_{1},X_{n+1}),
\]
组成的集合$\mathfrak{m}:=\{m_{n}\}_{n\ge 1}$，它确定$\ainf$结构。

把态射的次数整体平移一位的操作（悬垂）记作
\[
  s:\mathcal{C}(X,Y)\longrightarrow s\mathcal{C}(X,Y),\qquad
  (s\mathcal{C})^{i}(X,Y):=\mathcal{C}^{i+1}(X,Y),
\]
由此诱导的$\ainf$结构$sm_{n}(s^{-1})^{\otimes n}$仍用同一符号$m_{n}$表示。
此时这些次数与$n$无关，全为$1$。$\ainf$结构的条件式成为
\begin{equation}
\begin{aligned}
  0={}&\sum_{k+l=n+1}\sum_{j=0}^{k-1}
  (-1)^{|a_{12}|+\cdots+|a_{j(j+1)}|}\,
  m_{k}\bigl(a_{12},\dots,a_{j(j+1)},\\
  &m_{l}(a_{(j+1)(j+2)},\dots,a_{(j+l)(j+l+1)}),
  a_{(j+l+1)(j+l+2)},\dots,a_{n(n+1)}\bigr),
  \qquad(n=1,2,3,\dots)
\end{aligned}
\label{eq:c01-2}
\end{equation}
其中$a_{i(i+1)}\in s\mathcal{C}(X_{i},X_{i+1})$是次数为$|a_{i(i+1)}|$的态射。
换言之，
\[
  a_{i(i+1)}\in(s\mathcal{C})^{|a_{i(i+1)}|}(X_{i},X_{i+1})
  =\mathcal{C}^{|a_{i(i+1)}|+1}(X_{i},X_{i+1}).
\]

\transnote{这里的$m_{n}$与本节开头不是同一套次数约定：未悬垂时$\deg m_{n}=2-n$，
经$s$悬垂后$\deg m_{n}=1$，与$n$无关。式~\eqref{eq:c01-2}指数上的
$|a_{12}|+\cdots+|a_{j(j+1)}|$是悬垂空间$s\mathcal{C}$中的次数之和。混用两套约定
会把符号因子算错，这是本篇最易失手之处。}

$\ainf$范畴$\mathcal{C}$若没有高次积，即$m_{3}=m_{4}=\cdots=0$，就称为 DG 范畴。

\begin{definition}
  $\ainf$范畴$\mathcal{C}$若满足下述条件，则称为具有恒等态射（unital）：
  对每个对象$X\in\Ob(\mathcal{C})$，存在次数为零的态射
  $1_{X}\in\mathcal{C}^{0}(X,X)$，使对任意态射$a_{12},a_{23},\dots$
  （$a_{i(i+1)}\in\mathcal{C}(X_{i},X_{i+1})$）都有
  \[
    m_{2}(1_{X_{1}},a_{12})=a_{12}=m(a_{12},1_{X_{2}}),
  \]
  \[
    m_{n}(a_{12},\dots,a_{(i-1)i},1_{X_{i}},a_{i(i+1)},\dots,a_{n(n+1)})=0
    \quad(n\ge 3).
  \]
  \label{def:c01-2.1}
\end{definition}

\transnote{定义中第二式的右端，原文印作$m(a_{12},1_{X_{2}})$，脱漏下标$2$；
按前一式与$\ainf$结构应为$m_{2}(a_{12},1_{X_{2}})$。译文照原文保留$m$，
以便与底本逐字对照。}

% ===== 原书 p.45 =====
$\ainf$范畴的恒等态射还有若干条件更弱的概念；与它们区别时，满足上述条件的
$\ainf$范畴称为严格 unital（strictly unital）的。本文以下所说的$\ainf$范畴，
都指上述意义下具有恒等态射的$\ainf$范畴。

对任意$X,Y\in\Ob(\mathcal{C})$，上同调$H(\mathcal{C}(X,Y),m_{1})$上由$m_{2}$
诱导出态射的（结合的）合成。由此得到的$\mathbb{Z}$分次范畴记作$H(\mathcal{C})$。
特别地，零次上同调对积$m_{2}$封闭；把取零次上同调所得的范畴记作
$H^{0}(\mathcal{C})$。这里考虑的$\ainf$范畴$\mathcal{C}$都具有恒等态射，因此
$H(\mathcal{C})$、$H^{0}(\mathcal{C})$确实构成通常意义的范畴。

$\ainf$范畴之间的$\ainf$函子，定义为两个$\ainf$代数之间$\ainf$映射的类似物。
也就是说，$\ainf$函子$\mathfrak{f}:\mathcal{C}\to\mathcal{C}'$是指：态射
$f:\Ob(\mathcal{C})\to\Ob(\mathcal{C}')$，以及对任意
$X_{1},\dots,X_{k+1}\in\Ob(\mathcal{C})$确定的、次数为$(1-k)$的多重线性映射
\begin{equation}
  f_{k}:\mathcal{C}(X_{1},X_{2})\otimes\cdots\otimes\mathcal{C}(X_{k},X_{k+1})
  \longrightarrow\mathcal{C}'(f(X_{1}),f(X_{k+1}))
  \label{eq:c01-3}
\end{equation}
组成的集族，并且满足$\ainf$映射的条件。把它也按次数平移后的记法具体写出，就是
\[
\begin{aligned}
  \sum_{i\ge 1}\sum_{k_{1}+\cdots+k_{n}=n}
  &m'_{i}(f_{k_{1}}\otimes\cdots\otimes f_{k_{i}})(a_{12},\dots,a_{n(n+1)})\\
  &=\sum_{\substack{i+l+j=k\\ i+l+j=n}}
  f_{k}(1^{\otimes i}\otimes m_{l}\otimes 1^{\otimes j})(a_{12},\dots,a_{n(n+1)}),
  \qquad n=1,2,3,\dots
\end{aligned}
\]
其中对每个$i=1,\dots,n-1$，有$a_{i(i+1)}\in s\mathcal{C}(X_{i},X_{i+1})$。
也就是说，这里的$f_{k}$与$m_{k}$的情形相同，是由式~\eqref{eq:c01-3}的$f_{k}$
把次数平移一位诱导而来；因此与$k$无关，次数全为零。

\transnote{左端求和指标原文印作$k_{1}+\cdots+k_{n}=n$，但左端张量积
$f_{k_{1}}\otimes\cdots\otimes f_{k_{i}}$只有$i$个因子，指标应为
$k_{1}+\cdots+k_{i}=n$。右端两行下标$i+l+j=k$与$i+l+j=n$合起来，等于要求
$k=i+l+j=n$；常见写法只写这一条件的一行。两处均照原文排印。}

$n=1$时的条件表明，$f_{1}$给出从复形$(\mathcal{C}(X_{1},X_{2}),m_{1})$到
复形$(\mathcal{C}(X_{1},X_{2}),m_{1})$的链映射，这一点需要注意。

\transnote{原文这里把$f_{1}$的目标复形也印成$(\mathcal{C}(X_{1},X_{2}),m_{1})$，
与源复形完全相同；按式~\eqref{eq:c01-3}，目标复形应为
$(\mathcal{C}'(f(X_{1}),f(X_{2})),m_{1})$。译文照原文保留。}

进而，$n=2$时的条件表明$(f,f_{1})$给出从$H(\mathcal{C})$到$H(\mathcal{C}')$的
函子。把它记作$H(\mathfrak{f}):H(\mathcal{C})\to H(\mathcal{C}')$。

作为$\ainf$代数之间$\ainf$拟同构映射的推广，称满足下述条件的$\ainf$函子
$\mathfrak{f}=\{f;f_{1},f_{2},f_{3},\dots\}$为$\ainf$拟同构函子：
$f:\mathcal{C}\to\mathcal{C}'$为双射，且$f_{1}$给出复形之间的拟同构映射。特别地，
当$f_{1}$是复形之间的同构映射时，称$\mathfrak{f}$为$\ainf$同构函子。

$\ainf$范畴之间的自然等价关系，应当是通常意义下范畴等价的概念与$\ainf$拟同构
函子适当配合的产物。由于$H(\mathcal{C})$构成通常意义的范畴，按下面这样定义，
它实际上就是合适的概念。

\begin{definition}
  $\ainf$函子$\mathfrak{f}:\mathcal{C}\to\mathcal{C}'$称为$\ainf$等价函子，若
  $H(\mathfrak{f}):H(\mathcal{C})\to H(\mathcal{C}')$是范畴等价。若存在$\ainf$等价
  函子$\mathfrak{f}:\mathcal{C}\to\mathcal{C}'$，则称$\mathcal{C}$与$\mathcal{C}'$
  为$\ainf$等价的。
  \label{def:c01-2.2}
\end{definition}

由这个定义可知，$\ainf$拟同构函子是$\ainf$等价函子。

% ===== 原书 p.46 =====
给定$\ainf$代数时，必存在与其$\ainf$拟同型的极小$\ainf$代数[6]。$\ainf$范畴的
情形也完全一样，以下结论成立。

\begin{theorem}[{$\ainf$范畴的极小模型定理}]
  给定$\ainf$范畴$\mathcal{C}$时，存在某个极小$\ainf$范畴$\mathcal{C}'$与
  $\ainf$拟同构函子$\mathfrak{f}:\mathcal{C}'\to\mathcal{C}$。当$\mathcal{C}$
  具有恒等态射时，可取$\mathcal{C}'$也具有恒等态射，且$H(\mathfrak{f})$为范畴同构。
  \label{thm:c01-2.3}
\end{theorem}

$\mathcal{C}'$与$\mathfrak{f}$的具体构造法也存在（cf.\ [15, 7]）。事实上，把
$\ainf$代数情形已知的具体构造法，通过将$\mathcal{C}$视为$\ainf$代数来应用，
所得$\mathcal{C}'$上的$\ainf$结构确定$\ainf$范畴的结构（在无法合成的态射组上
取平凡值），这一点是显然的；同样，$\ainf$拟同构映射也给出$\ainf$拟同构函子
$\mathcal{C}'\to\mathcal{C}$。

已知不仅$\ainf$拟同构、$\ainf$等价函子也（正如其名）给出$\ainf$范畴之间的等价
关系（cf.\ [20]）。因此，与上述极小模型定理合起来，对任意$\ainf$范畴
$\mathcal{C}$，都可取到与其$\ainf$等价的极小$\ainf$范畴$\mathcal{C}'$，使任意
$X,Y\in\mathcal{C}'$在$H(\mathcal{C}')$中同构仅当$X=Y$。可以说，这就是
$\ainf$范畴的$\ainf$等价类作为代表元（representative）的“最简明”者。
